% ============================================================================
%  section_ikk_kernel_bridge.tex
%
%  IKK Theory to VSIM Kernel Bridge
%  Mapping the Inverse Kaluza-Klein Scale-Symmetry Framework
%  onto the VSEPR-SIM IDENTITY layer
%
%  Theory basis:  IKK_Doc_I_Schrodinger_Alternative.tex  (Doc I)
%                 IKK_Doc_II_Initial_Draft.tex            (Doc II)
%                 IKK_Doc_III_Kernel_V4_Projection.tex    (Doc III)
%
%  Code targets:  include/IDENTITY/identity_matrix.hpp
%                 include/IDENTITY/particle_identity.hpp
%                 include/IDENTITY/interaction_channel.hpp
%                 include/IDENTITY/interaction_candidate.hpp
%
%  Compile:  pdflatex section_ikk_kernel_bridge.tex  (run twice for TOC)
%
%  v5.2.0  |  WO-SM-IDENTITY  |  v5.0.0-main
% ============================================================================

\documentclass[11pt, a4paper, oneside]{article}

\usepackage[top=1.0in, bottom=1.0in, left=1.15in, right=1.15in,
			headheight=14pt]{geometry}
\usepackage[T1]{fontenc}
\usepackage{lmodern}
\usepackage{microtype}
\usepackage{amsmath, amssymb, amsthm, mathtools}
\usepackage{bm}
\usepackage{physics}
\usepackage{booktabs}
\usepackage{array}
\usepackage{tabularx}
\usepackage{listings}
\usepackage{xcolor}
\usepackage{enumitem}
\setlist{nosep, leftmargin=1.5em}
\usepackage{fancyhdr}
\usepackage{titlesec}
\usepackage[colorlinks=true, linkcolor=black, citecolor=black,
			urlcolor=blue]{hyperref}

% --- Code style --------------------------------------------------------------
\definecolor{codegray}{rgb}{0.5,0.5,0.5}
\definecolor{codeblue}{rgb}{0.13,0.13,0.9}
\definecolor{codegreen}{rgb}{0.0,0.5,0.0}
\definecolor{codebg}{rgb}{0.97,0.97,0.97}
\lstset{
  backgroundcolor=\color{codebg},
  basicstyle=\ttfamily\small,
  keywordstyle=\color{codeblue}\bfseries,
  commentstyle=\color{codegreen}\itshape,
  numberstyle=\tiny\color{codegray},
  numbers=left, numbersep=6pt,
  breaklines=true, frame=single,
  rulecolor=\color{codegray},
  captionpos=b, showstringspaces=false,
  tabsize=2,
  language=C++
}

% --- Math operators ----------------------------------------------------------
\newcommand{\Id}{\mathcal{I}}
\newcommand{\Anch}{\mathcal{A}}
\newcommand{\Proj}{P}
\newcommand{\Res}{\mathcal{R}}
\newcommand{\Pack}{\Pi}
\newcommand{\Scale}{S}

% --- Page style --------------------------------------------------------------
\pagestyle{fancy}
\fancyhf{}
\fancyhead[L]{\small\textit{IKK $\to$ VSIM Kernel Bridge}}
\fancyhead[R]{\small\textit{VSEPR-SIM v5.2.0 / IDENTITY layer}}
\fancyfoot[C]{\thepage}
\renewcommand{\headrulewidth}{0.4pt}

\titleformat{\section}{\Large\bfseries}{\S\thesection}{1em}{}
\titleformat{\subsection}{\large\bfseries}{\thesubsection}{1em}{}
\titleformat{\subsubsection}{\normalsize\bfseries}{\thesubsubsection}{1em}{}

\theoremstyle{definition}
\newtheorem{definition}{Definition}[section]
\theoremstyle{remark}
\newtheorem{remark}{Remark}[section]

% ============================================================================
\begin{document}

\begin{titlepage}
  \centering
  \vspace*{1.5cm}
  {\Huge\bfseries IKK Theory $\to$ VSIM Kernel Bridge\par}
  \vspace{0.8cm}
  {\Large Mapping Inverse Kaluza-Klein Scale Symmetry\\[0.3em]
  onto the VSEPR-SIM Identity Layer\par}
  \vspace{1.5cm}
  {\normalsize
	This document is the formal bridge between the Inverse Kaluza-Klein (IKK)
	theoretical framework (Documents I--III) and the \texttt{IDENTITY/} header
	library in VSEPR-SIM v5.2.0.  Each theoretical construct is mapped to its
	direct C++ counterpart, with design rationale and the specific open variables
	that remain formally undetermined.
  \par}
  \vspace{1.5cm}
  {\large VSEPR-SIM Theoretical Division\par}
  \vspace{0.3cm}
  {\large v5.2.0 --- WO-SM-IDENTITY\par}
\end{titlepage}

\tableofcontents
\newpage

% ============================================================================
\section{Overview and Purpose}
% ============================================================================

The Inverse Kaluza-Klein (IKK) framework (Doc I) proposes that the standard
quantum wavefunction $\Psi(\mathbf{r}, t)$ is a \emph{scale-projected shadow}
of a fuller identity structure $\Id$ defined over a scale ladder
$\{S_0, S_1, \ldots, S_D\}$.  The observed 3D spatial behaviour is the
$S_3$-level projection; corrections at higher scales produce residuals that
connect to dark-sector and quantum-gravity phenomenology.

VSEPR-SIM does not implement a wavefunction.  It is a deterministic atomistic
simulation engine.  The IKK framework provides three tools directly applicable
to the VSIM kernel at the subatomic annotation layer:

\begin{enumerate}
  \item \textbf{Scale ladder} --- a principled way to label which interactions
	are active for each particle type (IKK Doc I \S2, Doc II \S B).
  \item \textbf{Hidden-active condition} --- detection of particles whose
	net-zero observable channel conceals nonzero internal structure
	(IKK Doc I \S3.3, Doc II \S D).
  \item \textbf{Identity anchor and fuzz} --- the weighted center of an
	identity bundle and its spatial spread (IKK Doc I \S2.2--2.3).
\end{enumerate}

The SM-channel identity layer (\texttt{include/IDENTITY/}) is the correct
integration point for all three.  The Bridge65 gate remains in force:
subatomic energy-regime activation requires its own work order.  The IKK
additions are \emph{annotation enrichments} only --- they do not modify
classical MD positions or forces.

% ============================================================================
\section{Scale Ladder Mapping}
% ============================================================================

\subsection{IKK Scale Ladder (Doc I \S2.1, Doc II \S B)}

\begin{table}[h]
  \centering
  \begin{tabularx}{\linewidth}{c l l l l}
	\toprule
	$S_n$ & IKK DOF & Mediator & Obs.\ limit & VSIM interpretation \\
	\midrule
	$S_0$ & Binary existence $E\in\{0,1\}$ & --- & Planck & Particle exists in plasma \\
	$S_1$ & Linear vector & --- & Sub-nuclear & Reserved (not yet used) \\
	$S_2$ & Relational / angular & Gluon & Nuclear ($\sim$1 fm) & Color charge, Cornell force \\
	$S_3$ & Spatial / binding & Photon & Atomic ($\sim$1 \AA) & EM channel, Coulomb \\
	$S_4$ & Time-ordered & $W, Z$ bosons & $\sim10^{-18}$ m & Weak channel (reserved) \\
	$S_5$ & Mass-energy curvature & Graviton & Cosmological & Not modeled \\
	$S_D$ & Dark-scale & $X_D$ (unknown) & $> L_\text{EM}$ & Not modeled \\
	\bottomrule
  \end{tabularx}
  \caption{IKK scale ladder mapped to VSIM particle annotations.}
  \label{tab:scale-ladder}
\end{table}

\subsection{C++ Implementation: \texttt{ScaleLevel} and \texttt{scale\_mask}}

\begin{lstlisting}[caption={ScaleLevel enum and helpers in identity\_matrix.hpp}]
enum class ScaleLevel : uint8_t {
	S0 = 0,  // existence
	S1 = 1,  // vector (reserved)
	S2 = 2,  // color/gluon
	S3 = 3,  // EM/spatial
	S4 = 4,  // weak/time-ordered
	S5 = 5,  // curvature (not modeled)
	SD = 6,  // dark (not modeled)
};

inline constexpr uint8_t scale_bit(ScaleLevel s) noexcept { ... }
inline bool scale_active(uint8_t mask, ScaleLevel s) noexcept { ... }
\end{lstlisting}

Each \texttt{ParticleIdentity} carries a \texttt{scale\_mask} bitmask.  The
convenience constructors set defaults reflecting physics:

\begin{center}
  \begin{tabular}{l l}
	\toprule
	Particle archetype & \texttt{scale\_mask} bits \\
	\midrule
	Electron / positron & $S_0 + S_3 + S_4$ (EM + weak) \\
	Quark               & $S_0 + S_2 + S_3$ (color + EM) \\
	\bottomrule
  \end{tabular}
\end{center}

% ============================================================================
\section{Identity Bundle and Anchor}
% ============================================================================

\subsection{IKK Theory (Doc I \S2.2)}

A particle is an identity bundle:
\begin{equation}
  \Id_i(t) = \{c_{i1}, c_{i2}, \ldots, c_{im}\}
\end{equation}
with identity anchor:
\begin{equation}
  \Anch_i(t)
  = \frac{\sum_k \alpha_{ik}\, c_{ik}(t)}{\sum_k \alpha_{ik}}
  \label{eq:anchor}
\end{equation}
and fuzz radius:
\begin{equation}
  r_f^2
  = \frac{\sum_k \alpha_k \left\|P_{3D}(c_k) - \vec{r}^{\,\text{obs}}\right\|^2}
		 {\sum_k \alpha_k}
  \label{eq:fuzz}
\end{equation}

The fuzz ratio:
\begin{equation}
  \Phi_S = \frac{r_f}{\ell_S}
  \qquad
  \Phi_S \ll 1 \implies \text{point-like at scale } S
\end{equation}

\subsection{VSIM Mapping}

VSIM particles are point-like at $S_3$ (classical MD positions are exact).
The \texttt{r\_fuzz} field on \texttt{ParticleIdentity} stores a caller-updated
dimensionless spread proxy that can be populated when the particle is
participating in a QGP substep or a pair interaction with significant kinematic
spread.  The \texttt{fuzz\_ratio()} helper in \texttt{identity\_matrix.hpp}
converts this to the IKK $\Phi_S$ value.

\begin{lstlisting}[caption={fuzz\_ratio() in identity\_matrix.hpp}]
// Phi_S = r_fuzz / ell_S
inline double fuzz_ratio(double r_fuzz, ScaleLevel s,
						 double ell_S3_ang = 1.0,
						 double ell_S2_fm  = 1.0e-5);
\end{lstlisting}

\begin{remark}
  The IKK identity anchor~(\ref{eq:anchor}) is a weighted mean of multi-channel
  component values, not a simple 3D position.  In VSIM the 3D position is held
  by the classical MD layer (\texttt{QCDParticle::pos}); the identity anchor is
  the abstract center of the \emph{annotation} layer and is not directly
  observable.  The two are related by $\vec{r}^{\,\text{obs}} = P_{3D}(\Anch_i)$.
\end{remark}

% ============================================================================
\section{Hidden-Active Condition}
% ============================================================================

\subsection{Theory (Doc I \S3.3, Doc II \S D)}

For a component vector $\mathbf{c}_j \in \mathbb{R}^m$ with projection
weight $\bm{\alpha}$ and weight matrix $\mathbf{W} \succeq 0$:
\begin{align}
  I_j &= \bm{\alpha}^\top \mathbf{c}_j \label{eq:observable} \\
  H_j &= \mathbf{c}_j^\top \mathbf{W}\, \mathbf{c}_j \label{eq:hidden}
\end{align}

The \textbf{hidden-active condition}:
\begin{equation}
  \boxed{I_j = 0, \quad H_j > 0}
  \label{eq:hidden-active}
\end{equation}

\textbf{Canonical example: neutron-like charge balance.}\\
A neutral baryon composed of (up, down, down) quarks has:
\begin{align}
  I_\text{charge}
	&= +\tfrac{2}{3} - \tfrac{1}{3} - \tfrac{1}{3} = 0 \label{eq:neutron-Ij} \\
  H_\text{charge}
	&= \left(\tfrac{2}{3}\right)^2
	 + \left(\tfrac{1}{3}\right)^2
	 + \left(\tfrac{1}{3}\right)^2
	 = \tfrac{2}{3} \neq 0 \label{eq:neutron-Hj}
\end{align}

The particle is charge-neutral at the $K_\text{EM}$ projection level but
carries nonzero internal charge intensity visible to higher-scale projections.
This is not an approximation; it is a structural consequence of packing.

\subsection{VSIM Implementation}

Three hidden-intensity channels are tracked per \texttt{ParticleIdentity}:

\begin{lstlisting}[caption={Hidden intensity fields in particle\_identity.hpp}]
double H_charge {0.0};   // hidden charge intensity  H = sum q_k^2
double H_color  {0.0};   // hidden color intensity
double H_spin   {0.0};   // hidden spin intensity
\end{lstlisting}

The method \texttt{compute\_hidden\_intensity()} populates all three using the
single-particle or quark-triplet proxy:

\begin{lstlisting}[caption={compute\_hidden\_intensity() key logic}]
// H_charge for elementary particle: H = q^2
// H_charge for baryon proxy (isospin bookkeeping):
//   n_up * (2/3)^2 + n_down * (1/3)^2
void compute_hidden_intensity() noexcept;

// Query: net-zero charge but H_charge > 0?
bool hidden_active_charge(double tol = 1e-9) const noexcept;
\end{lstlisting}

\texttt{evaluate\_pair()} in \texttt{interaction\_candidate.hpp} populates the
hidden-active flags for each member of the pair:

\begin{center}
  \begin{tabular}{l l}
	\toprule
	Flag & Meaning \\
	\midrule
	\texttt{a\_hidden\_active\_charge} & particle $a$: $I_\text{charge}=0$, $H_\text{charge}>0$ \\
	\texttt{b\_hidden\_active\_charge} & particle $b$ \\
	\texttt{a\_hidden\_active\_color}  & particle $a$: color-neutral but $H_\text{color}>0$ \\
	\texttt{b\_hidden\_active\_color}  & particle $b$ \\
	\texttt{pair\_has\_hidden\_charge\_activity} & logical OR of the charge flags \\
	\bottomrule
  \end{tabular}
\end{center}

These flags are available to the live event logger and to any future kernel
that needs to distinguish a ``truly neutral'' particle from a
``net-zero-but-internally-structured'' one.

% ============================================================================
\section{Matrix Packing Chain}
% ============================================================================

\subsection{IKK Theory (Doc II \S C)}

The quark component matrix:
\begin{equation}
  \mathbf{C}_q =
  \begin{bmatrix}
	q_1 & c_1 & \eta_1 \\
	q_2 & c_2 & \eta_2 \\
	q_3 & c_3 & \eta_3
  \end{bmatrix}
  \in \mathbb{R}^{3\times 3}
  \qquad \text{(charge, color, helicity columns)}
\end{equation}

packed to particle identity row-vector:
\begin{equation}
  \mathbf{I}_p = \Pack_q(\mathbf{C}_q) \in \mathbb{R}^{1\times 3}
  \qquad [Q_p,\; M_p,\; S_p]
\end{equation}

The recursive packing ladder (Doc II eq.\ C.4):
\begin{equation}
  \underbrace{\text{quarks}}_{S_2}
  \xrightarrow{\Pack_2}
  \underbrace{\text{particles}}_{S_3}
  \xrightarrow{\Pack_3}
  \underbrace{\text{atoms}}_{S_3}
  \xrightarrow{\Pack_3}
  \underbrace{\text{molecules}}_{S_3}
  \xrightarrow{\Pack_4}
  \underbrace{\text{materials}}_{S_4}
\end{equation}

\subsection{VSIM Mapping}

The $3\times 3$ component matrix $\mathbf{C}_q$ corresponds to the
\texttt{Z3} field of \texttt{ParticleIdentity}.  For quarks, the three rows
hold one quark each; for mesons or baryons, the whole structure represents
the internal composition.  At present \texttt{Z3} is initialised to the
identity matrix; the packing operator $\Pack_q$ is a future kernel.

The entropy and complexity measures (Doc II \S C.3--C.4):
\begin{align}
  \mathcal{S}_n &= k_B \ln \Omega(\mathbf{I}^{(n)}) \\
  \mathcal{C}_n &= f(\mathbf{I}^{(n+1)},\, \Res_n)
\end{align}
are not yet computed; they appear as open variables in the kernel
development roadmap.

% ============================================================================
\section{Hidden Intensity: $Z_2$ and $Z_3$ as Inner-Structure Proxies}
% ============================================================================

The two $Z$-matrices in \texttt{ParticleIdentity} encode the particle's
field configuration at two resolutions:

\begin{center}
  \begin{tabular}{l l l}
	\toprule
	Matrix & IKK analog & Current role \\
	\midrule
	\texttt{Z2} ($2\times 2$) & Runtime identity state & kappa\_Z (2x2 opposition) \\
	\texttt{Z3} ($3\times 3$) & Component packing matrix $\mathbf{C}_q$ & Future packing kernel \\
	\bottomrule
  \end{tabular}
\end{center}

The \texttt{opposition()} function on \texttt{Z2} matrices produces
\texttt{kappa\_Z}, the leading term in the weighted kappa\_total:
\begin{equation}
  \kappa_Z = 1 - \frac{\langle Z_2^{(a)}, Z_2^{(b)}\rangle_F}
					  {\|Z_2^{(a)}\|_F \|Z_2^{(b)}\|_F + \varepsilon}
\end{equation}
This is the most direct numerical realisation of the IKK hidden-structure
opposition concept at the 2D identity level.

The weight $w_Z = 0.40$ (40\% of kappa\_total) reflects the design choice
that internal-structure opposition is the dominant interaction discriminator.

% ============================================================================
\section{Dark-Scale Residual and Open Variables}
% ============================================================================

The IKK dark-sector equations (Doc I \S6, Doc II \S E) define:
\begin{equation}
  \rho_\text{dark}(x,y,z;\,t)
  = \int \rho(x,y,z,w;\,t)
	\bigl[K_G(w) - K_\text{EM}(w)\bigr]\, dw
\end{equation}

\textbf{VSIM status:} The dark coordinate $w$ and both projection kernels
$K_G$, $K_\text{EM}$ are open variables (Doc II open-variable register
\S J).  The \texttt{scale\_mask} field with bit $S_D = 6$ is the forward
compatibility hook; no dark-scale force or density computation is active.

A reference C++ implementation of \texttt{dark\_density\_residual()} using
trapezoidal integration over a $w$-grid appears in Doc II \S I (C++ code
section) and in the \texttt{ikk::} namespace sketch distributed alongside
that document.

% ============================================================================
\section{Projection Failure and State Loss}
% ============================================================================

IKK Doc I \S4 establishes that when projection $P_S$ is non-injective,
exact reconstruction is impossible.  The VSIM \texttt{I\_recoverable} field
tracks the recoverable fraction of the identity state after each interaction:

\begin{align}
  \Delta I   &= I_\text{out} - I_\text{in}   \quad\text{(signed change)} \\
  L_{ij}     &= -\Delta I                    \quad\text{(state loss, } \geq 0\text{)}
\end{align}

State loss $L_{ij} > 0$ corresponds to IKK \emph{destructive inverse packing}:
information available at the higher scale ($\Id$) is not recovered at the
lower projection ($\mathcal{O}_{S_3}$).  The annihilation gate requires
$L_{ij} > L_\text{ann}$ (default 0.50) to confirm a full annihilation event.

\begin{remark}
  The IKK framework connects this to the quantum no-cloning theorem and
  to the irreversibility of measurement.  In VSIM this is operationalised
  as an interaction-record bookkeeping quantity, not a wavefunction collapse.
\end{remark}

% ============================================================================
\section{Relation-Particle and Entanglement}
% ============================================================================

IKK Doc I \S5 / Doc II \S G defines the relation-particle object:
\begin{equation}
  \mathcal{R}_{AB}
  = \bigl[0,\; 0,\; 0,\; 0 \mid \mathbf{T}_{AB}\bigr]
  \qquad
  \mathbf{T}_{AB} \in \{-1,\, 0,\, +1\}^m
\end{equation}

$\mathcal{R}_{AB}$ is a real, deterministic, non-spatiotemporal structure.
Measurement at $A$ and at $B$ both read the same relation; no signal travels.

\textbf{VSIM status:} The trinary vector $\mathbf{T}_{AB}$ and its dynamics
are open variables.  The \texttt{InteractionCandidate} struct is the natural
carrier for a relation-particle record when this is implemented; the
\texttt{pair\_has\_hidden\_charge\_activity} flag is the first step toward
identifying pairs that would carry a non-trivial $\mathbf{T}_{AB}$.

% ============================================================================
\section{Effective Proper Time}
% ============================================================================

The IKK extended proper-time element (Doc I \S7, Doc II \S F):
\begin{equation}
  d\tau_\text{eff}
  = dt\sqrt{1 - \frac{v^2}{c^2} + \frac{2\Phi_G}{c^2} + \chi_w(w)}
\end{equation}

reduces to standard SR proper time when $\chi_w = 0$.

\textbf{VSIM status:} VSIM is a classical MD engine.  No relativistic proper
time is computed.  The $T_A|B$ dual-timescale architecture (WO-63A) provides
an operational time-dilation analog for QCD substeps ($\tau_\text{scale}$),
but this is a scale-bridging approximation, not the IKK dark-scale correction.
The \texttt{scale\_mask} bit $S_4$ (weak / time-ordered) is the forward
hook for proper-time corrections if the engine is ever extended to
relativistic integration.

% ============================================================================
\section{Code Reference: Files and Namespaces}
% ============================================================================

\begin{center}
  \begin{tabularx}{\linewidth}{l X}
	\toprule
	File & IKK concept implemented \\
	\midrule
	\texttt{include/IDENTITY/identity\_matrix.hpp}
	  & \texttt{ScaleLevel}, \texttt{scale\_bit/active}, \texttt{fuzz\_ratio()},
		\texttt{hidden\_intensity\_diagonal()}, \texttt{quark\_charge\_hidden\_intensity()} \\
	\texttt{include/IDENTITY/particle\_identity.hpp}
	  & \texttt{ParticleIdentity}: \texttt{scale\_mask}, \texttt{r\_fuzz},
		\texttt{H\_charge/color/spin}, \texttt{compute\_hidden\_intensity()},
		\texttt{hidden\_active\_charge()} \\
	\texttt{include/IDENTITY/interaction\_channel.hpp}
	  & Conservation gates $G_E \cdots G_s$, kappa scoring,
		nine-channel row \\
	\texttt{include/IDENTITY/interaction\_candidate.hpp}
	  & \texttt{evaluate\_pair()}: hidden-active flags
		\texttt{a/b\_hidden\_active\_charge/color},
		\texttt{pair\_has\_hidden\_charge\_activity} \\
	\texttt{include/IDENTITY/identity.hpp}
	  & Umbrella header \\
	\texttt{include/COLLISION/qcd\_force.hpp}
	  & \texttt{qcd\_accumulate\_forces\_sm()}: SM-gated Cornell wrapper \\
	\texttt{include/IKK/kernel\_v4.hpp}
	  & \texttt{ikk::KernelV4}: two-channel bilinear projection (Doc III);
		\texttt{ikk::KernelV4Multi<N\_CHAN,N\_STATE>}: generalised engine;
		\texttt{ikk::mixing\_matrix\_2ch()}: standard $\mathbf{M}_2$ \\
	\bottomrule
  \end{tabularx}
\end{center}

All \texttt{IDENTITY/} code is in the \texttt{vsepr::identity} namespace.
The bilinear projection kernel lives in the \texttt{ikk::} namespace
(\texttt{include/IKK/kernel\_v4.hpp}), separate from the identity layer,
to keep the radial-density engine independent of particle bookkeeping.

% ============================================================================
\section{Open Variables Register (VSIM perspective)}
% ============================================================================

\begin{table}[h]
  \centering
  \begin{tabularx}{\linewidth}{l l X}
	\toprule
	Variable & IKK source & VSIM status \\
	\midrule
	$\chi_w(w)$ & Doc I \S7 / Doc II F.5
	  & Dark-scale metric correction. Not modeled. \texttt{scale\_mask} $S_D$ is the hook. \\
	$K_G(w),\, K_\text{EM}(w)$ & Doc I \S6 / Doc II E.3
	  & Projection kernels. C++ skeleton in Doc II \S I. Not active. \\
	$\Pack_n$ & Doc II B.3
	  & Packing operator. \texttt{Z3} is the target matrix. Axiom only. \\
	$\mathbf{T}_{AB}$ & Doc II G.2
	  & Trinary relation vector. \texttt{InteractionCandidate} is the carrier. Not computed. \\
	$\mathbf{W}$ & Doc II D.3
	  & Hidden-intensity weight matrix. Diagonal $\mathbf{W}=\mathbf{I}$ assumed currently. \\
	$\hat{H}_w$ & Doc I \S3.2
	  & Dark-scale Hamiltonian. No quantitative predictions without this. \\
	$\mathcal{S}_n,\, \mathcal{C}_n$ & Doc II C.3--C.4
	  & Packing entropy and complexity. Roadmap item. \\
	\bottomrule
  \end{tabularx}
  \caption{Open variables from IKK theory and their VSIM status.}
  \label{tab:open-vars}
\end{table}

% ============================================================================
\section{Summary}
% ============================================================================

The IKK scale-symmetry framework provides the VSEPR-SIM IDENTITY layer with:

\begin{enumerate}
  \item A principled scale-ladder labeling system (\texttt{ScaleLevel} enum,
	\texttt{scale\_mask} bitmask) that connects VSIM particle types to their
	active force channels.
  \item The hidden-active condition ($I_j=0$, $H_j>0$) implemented as
	\texttt{H\_charge/color/spin} fields, \texttt{compute\_hidden\_intensity()},
	and hidden-active flags on \texttt{InteractionCandidate}.
  \item The fuzz proxy (\texttt{r\_fuzz}) and \texttt{fuzz\_ratio()} helper
	that operationalise $\Phi_S = r_f / \ell_S$ for future spatial-spread tracking.
  \item The quark-component-matrix slot (\texttt{Z3}) as the target for the
	future packing operator $\Pack_2$.
  \item A documented open-variable register linking every undefined IKK quantity
	to the VSIM code location where it will eventually be implemented.
\end{enumerate}

The bridge does not modify the classical MD layer.  All additions are
annotation-only and live entirely within the \texttt{vsepr::identity}
namespace, respecting the Bridge65 separation doctrine.

% ============================================================================
\section*{Reading Path}
% ============================================================================

\begin{enumerate}
  \item \textbf{IKK Doc I} --- theoretical foundation and Schr\"{o}dinger alternative
  \item \textbf{IKK Doc II} --- scale axioms, matrix packing chain, hidden intensity,
		dark matter kernel equations, Python and C++ prototypes
  \item \textbf{IKK Doc III} --- Kernel V4 bilinear projection form
  \item \textbf{section\_sm\_identity\_layer} --- full SM-channel IDENTITY layer
		documentation (conservation gates, $\kappa_\text{total}$, $\Delta I$,
		live event types)
  \item \textbf{This document} --- IKK $\to$ VSIM kernel mapping
\end{enumerate}

\end{document}
